% 原文 pp.276–283（印刷 pp.1394–1401）
% 第 19 讲第 3 部分：四维 $\mathcal{N}=2$ 超对称 Yang--Mills 理论（三）——拓扑应用。
% 承接 lec19_p2.tex，故不再写 \section。
% 注意：原书 pp.1394–1401 上的一切展示公式均未编号，故按规范 §3 一律用 \[...\]，不设 \tag。

% ===== 原文 p.276 =====
（承上页）这里 $A$ 是线丛 $\det(S^+\otimes\mathcal{L})$ 上的一个 $U(1)$ 联络，而 $M$ 是 $S^+\otimes\mathcal{L}$ 的一个截面。当然，我们必须对所有复线丛求和。

实际上这个术语并不精确。我们非正式地写成的 $S^+\otimes\mathcal{L}$，在数学上通常称为 $X$ 上的一个 $\mathrm{Spin}^c$ 结构（$\mathrm{Spin}^c$ structure）。因此我们实际上是在对 $X$ 上一切 $\mathrm{Spin}^c$ 结构求和；这类结构的集合构成复线丛同构类群之下的一个主齐性空间（principal homogeneous space）。关于从大 $u$ 处的描述到 $u=\pm1$ 处的描述的对偶变换如何导致 $\mathrm{Spin}^c$ 结构，其详细解释见参考文献中所引的 Moore--Witten 论文。

% [Tutor Note: 关于「$\mathrm{Spin}^c$ 结构构成主齐性空间」这句话，值得把群作用写清楚：若 $\lambda$ 是一个 $\mathrm{Spin}^c$ 结构（即某个 $S^+\otimes\mathcal{L}$），则任何复线丛 $\mathcal{L}'$ 都通过 $\lambda\mapsto\lambda\otimes\mathcal{L}'$ 给出另一个 $\mathrm{Spin}^c$ 结构，而同构类群正是 $H^2(X;\mathbb{Z})$；此作用自由且传递，故 $\mathrm{Spin}^c$ 结构之集是 $H^2(X;\mathbb{Z})$ 的主齐性空间（不是群，因为没有典范的单位元）。这正是后文对所有 $\lambda$ 求和时不会重复计数、也不会漏掉任何一支的根据。另请注意 $\operatorname{rank}S^+=2$，故 $c_1(\det(S^+\otimes\mathcal{L}))=c_1(S^+)+2c_1(\mathcal{L})$，此式的「一半」正是 p.278 处把 $\lambda$ 的第一 Chern 类定义为一个半整数值的由来。]

$M$ 只在 $u=\pm1$ 处出现，这正是测度在这些点上有 delta 函数型贡献的原因。含 $M\ne0$ 的方程的解无法从 $u=\pm1$ 处有意义地延拓出去，因为 $M$ 会获得质量。我们可以换一种更形式化的方式来陈述这一论证：回顾 §19.2 的最末尾，我们在那里得知，把超多重态（hypermultiplet）计入之后，$Q$-不变场的条件之一是
\[
0=\phi(M).
\]
在目前的情形下，在 $u=1$ 附近的有效理论中，$\phi$ 就是光子的超伴侣（superpartner）$a$——这里所用的 Lagrange 量描述在 $u=1$ 附近是好的。超多重态 $M$ 在这种描述中具有电荷 $1$ 或 $-1$ 的分量；把这些分量记作 $T,\bar T$。于是 $\phi(M)$ 的分量为 $aT$ 与 $-a\bar T$，从而方程 $0=\phi(M)$ 化为
\[
0=aT=a\bar T.
\]
因此有两条互不相同的分支：其一为 $a=0$ 而 $T,\bar T$ 可以非零，这对应于支撑在 $u=1$ 处的 delta 函数；其二为 $T=\bar T=0$ 而 $a$ 可以为零，这对应于来自 $u$ 平面上一般点的贡献，此时超多重态被置为零。

初看之下，人们会以为这两类贡献之间未必有明确的分界，因为两条分支也许会在 $a=M$ 处相遇。然而请注意到：若 $a=M=0$，则剩下的方程化归为
\[
F^+=0,
\]
即 $U(1)$ 规范理论的瞬子（instanton）方程。若 $b_2^+(X)\ge1$，则对 $X$ 上一般的度规不存在 $U(1)$ 瞬子；而若 $b_2^+(X)>1$，则在一族一般的单参数度规族中不存在这样的解。\footnote{单参数度规族之所以重要，是因为：在证明某个用度规定义的表达式（例如 Donaldson 不变量或量子场论关联函数）实际上与度规无关时，人们通常考察在两张给定的度规之间插值的单参数度规族。} 因此，若 $b_2^+(X)>0$，两条分支的贡献是彼此良好分离的；并且当 $b_2^+(X)>1$ 时，这种分离与度规的选取无关。$b_2^+(X)=0$ 的情形还需要更多研究（目前尚未有人做过）。

% [Tutor Note: 「两条分支也许会在 $a=M$ 处相遇」一句在原文中写得过简：$a$ 是标量（光子超伴侣的期望值），$M$ 是超多重态场，两者不能相等；作者的意思是二者同时为零，紧接着的「若 $a=M=0$」才是精确说法。另，「$F^+=0$ 无解」的理由这里也跳过了：$F^+=0$ 意味着曲率是反自对偶（ASD）的调和二形式，于是 $c_1$ 必须落在实子空间 $H^2_-(X;\mathbb{R})$（维数 $b_2^-$）之中；而对一般度规，非零整系数上同调类不会恰好落在其中，故只有 $c_1=0$ 的解。注意 $c_1\in H^2_-$ 等价于 $b_2^+$ 个方程 $\langle\omega_i,c_1\rangle=0$（$\omega_i$ 为 $H^2_+$ 的基）：$b_2^+=1$ 时这只是单个方程，故在单参数度规族中一般会有解（度规空间中余维 $1$ 的「壁」），而 $b_2^+>1$ 时方程个数多于参数个数、一般无解——这正是原文「$b_2^+\ge1$ 时一般度规无解、$b_2^+>1$ 时一般单参数族亦无解」的确切含义，也是 Donaldson 理论在 $b_2^+(X)=1$ 时必须处理壁穿越（wall crossing）的几何根源。]

% [Tutor Note: 这条脚注所说的是拓扑场论中「度规无关性」的标准证明策略：设 $g_0,g_1$ 为两张度规，取道路 $g(t)$，若能证明 $\frac{d}{dt}\langle\mathcal{O}\rangle_{g(t)}=0$（其技术核心是 $Q$-正合性：$\frac{\partial}{\partial t}$ 作用在作用量上给出 $\{Q,\Lambda_t\}$ 形式的项，从而在 $Q$-闭可观测量的期望值中积掉），则两端相等。正文在 $b_2^+(X)>1$ 时援用「一般的单参数族中无解」，正是为了排除道路中途出现瞬子模空间改变拓扑（即壁穿越）的情形；注意 $b_2^+=1$ 时这个论证失效，必须另行讨论。]

\subsection{低能有效理论中的算符}
% 原文 §19.4.1：Operators in the effective low energy theory

为了通过低能有效场论计算关联函数，我们需要知道基本（即高能）理论中所定义的算符在低能有效理论中的像。在定义某个算符的低能有效版本时，我们必须把一切高能模式都积掉。这个操作与取乘积并不对易。因此，尽管我们将看到有些基本算符以显然的方式对应，乘积却会出些古怪的事情。

首先，在高能理论中我们有算符 $u=\operatorname{Tr}\phi^2$。我们已经看到，这个算符在低能极限下流向一个我们能显式写出的算符 $u$。这一陈述等价于说：精确的真空模空间（moduli space of vacua）就是 $u$ 平面，而我们取算符 $u$ 的期望值来标记各个真空。

% ===== 原文 p.277 =====
当然，高能理论中 $x$ 点处的局部算符映为低能理论中同一点 $x$ 处的局部算符。因此，高能理论中局部算符 $\mathcal{O}$ 的导数映为其像 $\bar{\mathcal{O}}$ 的导数；类似地，高能理论中的 $\int_\Sigma\mathcal{O}$ 流向 $\int_\Sigma\bar{\mathcal{O}}$。由于低能理论仍具有 $\mathcal{N}=2$ 超对称，它拥有算符 $\bar Q$ 与 $\bar K$，它们分别是高能理论中算符 $Q$ 与 $K$ 的像。于是 $K^nu$ 在低能理论中的像就是 $\bar K^nu$。因此，只要 $\Sigma_1,\Sigma_2$ 是 $X$ 中不相交的闭链，取乘积的操作就与对高能模式积分对易。这意味着：高能算符
\[
\int_{\Sigma_1}K^{n_1}u\cdot\int_{\Sigma_2}K^{n_2}u
\]
在低能理论中变成各个算符之像的乘积：
\[
\int_{\Sigma_1}\bar K^{n_1}u\cdot\int_{\Sigma_2}\bar K^{n_2}u,
\]
只要
\[
\Sigma_1\cap\Sigma_2=\emptyset.
\]
在过渡到低能有效理论时，乘积不被保持的场合正是闭链相交之时。此时答案的一般形式是：高能算符的乘积
\[
\int_{\Sigma_1}K^{n_1}u\cdot\int_{\Sigma_2}K^{n_2}u
\]
在低能极限下流向如下形式的算符：
\[
\int_{\Sigma_1}\bar K^{n_1}u\cdot\int_{\Sigma_2}\bar K^{n_2}u+\int_{\Sigma_1\cap\Sigma_2}\bar K^{n_1+n_2-4}f(u),
\]
其中 $f(u)$ 是某个函数。（若这两条闭链横截相交，则 $\Sigma_1\cap\Sigma_2$ 的维数为 $n_1+n_2-4$。）多出这一项的理由是：若把大质量自由度积掉给出 $\mathcal{O}_i(x)\mapsto\bar{\mathcal{O}}_i(x)$，那么对乘积一般会给出
\[
\mathcal{O}_1(x)\mathcal{O}_2(y)\mapsto\bar{\mathcal{O}}_1(x)\bar{\mathcal{O}}_2(y)+\delta(x-y)\tilde{\bar{\mathcal{O}}}_{12}(x)+\cdots
\]
其中 $\tilde{\bar{\mathcal{O}}}(x)$ 是低能理论中某个适当的算符（$\cdots$ 项正比于 $\delta(x-y)$ 的导数）。这个公式所断言的其实只是两件事：若 $x\ne y$，则积掉大质量自由度的操作可以分别作用于两个算符；而若不要求 $x\ne y$，则存在一个新的、支撑在 $x=y$ 处的分布型贡献，关于它的细节公式未作任何断言。换句话说，只有当 $x$ 与 $y$ 相遇时才可能发生新的事情，而除了「公式中可能出现新的局部算符」之外，我们对其间发生什么不作一般性的断言。在我们的情形中，若 $\mathcal{O}_i=K^{n_i}u$，则（由于 $Q$-不变性以及模去 $\{Q,\cdots\}$）必定有
\[
\tilde{\bar{\mathcal{O}}}_{12}=\bar K^{n_1+n_2-4}f(u),
\]
其中 $f(u)$ 是某个全纯函数。把这个结果对局部算符积分，就给出上面关于乘积
\[
\int_{\Sigma_1}K^{n_1}u\cdot\int_{\Sigma_2}K^{n_2}u
\]
在低能理论中的像的公式。

% [Tutor Note: 相交项中「$n_1+n_2-4$」这个指数是可以用形式次数一眼验证的，作者未写出：在本讲所用的扭变（twisting）理论中，$K$ 是一个把形式次数提高 $1$ 的超荷（作为微分形式算符，$\{K,\psi_\mu\}\sim\partial_\mu$），故 $K^nu$ 是一个 $n$ 次形式，$\int_{\Sigma_d}K^nu$ 只有在 $d=n$（$\Sigma_d$ 为 $d$ 维闭链）时才是零次、无标度依赖的拓扑可观测量的正确搭配（检验：$u$ 是 $0$ 形式、$\int_{\text{pt}}u$；$K^2u$ 是 $2$ 形式、$\int_{\Sigma_2}K^2u$，其低能展开 $(\partial u/\partial a)F+(\partial^2u/\partial a^2)\psi\wedge\psi$ 确实由 $2$ 形式组成）。两条闭链若横截相交，则 $\dim(\Sigma_1\cap\Sigma_2)=n_1+n_2-4$，而修正项的算符次数恰为 $n_1+n_2-4$，与公式自洽。指标易错之处在于：若把 $K$ 与 $\bar K$ 混用（低能理论中只保留了 $\bar K$ 的像，因为 $K,\bar K$ 在 $\mathcal{N}=2$ 中属于不同的 $\mathcal{N}=1$ 子代数），指数虽仍相同，但全纯性的结论 $f=f(u)$（而非 $\bar u$）就会丢掉依据。]

% ===== 原文 p.278 =====
若我们考虑三个或更多这样的算符（比如说在闭链 $\Sigma_1,\Sigma_2,\Sigma_3$ 上积分），则还会有更多同源的新项，来自三重交 $\Sigma_1\cap\Sigma_2\cap\Sigma_3$。

让我们假定 $b_1(X)=0$ 来简化讨论。这时我们只需考虑算符 $u,K^2u,K^4u$。但从 $K^4u$ 得不到多少新东西，因为它是数值不变量，而且恰好就是第二 Chern 类，于是把它插入关联函数只是把该丛上的期望值乘以瞬子数 $k$。从今以后我们就把这个算符搁置一旁。在剩下的算符之间不存在三重交，只要我们作如下假定：积分 $K^2u$ 所沿的那些二闭链处于一般位置（借助拓扑不变性，我们总是可以这样假定）。此外，唯一的双重交是有限个点，即一对二闭链横截相交之处。在这种情形下，高能算符的乘积
\[
\int_{\Sigma_1}K^2u\cdot\int_{\Sigma_2}K^2u
\]
在低能有效理论中的像将具有如下形式：
\[
\int_{\Sigma_1}\bar K^2u\cdot\int_{\Sigma_2}\bar K^2u+(\Sigma_1\cap\Sigma_2)T(u),
\]
其中 $T$ 是 $u$ 的全纯函数，而 $\Sigma_1\cap\Sigma_2$ 表示 $\Sigma_1$ 与 $\Sigma_2$ 的代数相交数。$T(u)$ 之所以全纯，理由是：由于 $Q$-闭算符映为 $Q$-闭算符，公式中的修正项必须是 $Q$-闭的；另一方面，模去 $Q$-正合项之后，任何 $Q$-闭的标量算符都是 $u$ 的全纯函数。

% [Tutor Note: 「$K^4u$ 恰好就是第二 Chern 类」这句话依赖一个归一化约定，作者没有写出：在扭变理论中 $K^4u$ 的低能展开为 $e''(a)F\wedge F+\psi\wedge(\cdots)$（见 p.280 第一个公式），而 $F\wedge F$ 的积分就是（$8\pi^2$ 倍）瞬子数 $k=\frac{1}{8\pi^2}\int_X\operatorname{Tr}F\wedge F=\frac{1}{8\pi^2}\int_X F\wedge F$（阿贝尔情形），故 $\int_X K^4u\!\propto\! k$。若采用另一种常见归一化 $\frac{1}{4\pi^2}\int F\wedge F$（例如把 $c_2$ 写成 $\frac{1}{4\pi^2}$ 倍），则本页之后所有 $2$ 的幂次与 $q^{\lambda^2}$ 中的 $\lambda^2$ 都会差一个因子 $2$。核对 p.279、p.281、p.282 的最终公式时，这是最容易出错的一处。]

\subsection{带超多重态的 $U(1)$ 理论的预备练习}
% 原文 §19.4.2：Warm-up with $U(1)$ theory with a hypermultiplet

现在我们要做一个练习性的计算，它有助于学会把 $SU(2)$ 理论与 $U(1)$ 理论联系起来。但我们要在一个更简单的理论里做这个练习，即带一个超多重态、且 Lagrange 量取「极小」形式的 $U(1)$ 理论；所谓极小，是指我们将省略某些耦合。我们干脆就取平直空间中、在 $\tau$ 为常数时写下的矢量多重态与超多重态的 Lagrange 量。令
\[
q=\exp(2\pi i\tau).
\]
李代数 $\mathfrak{u}(1)$ 唯一的不变多项式就是 $a$ 的多项式。我们假定 $H^1(X)=0$，于是所讨论的算符就是零维算符 $\mathcal{O}^{(0)}=a$ 与二维算符 $\mathcal{O}^{(2)}=c_1(\mathcal{L})$（其中 $\mathcal{L}$ 是规范场作为联络所在的那个线丛）。同样地，后一个算符的插入只是把关联函数乘以该线丛的第一 Chern 类，因而并无新意。也就是说，我们需要计算其期望值的算符，只是全纯（或多项式）函数 $P(a)$。

回忆一下，在这个理论中，超多重态场 $M$ 的玻色部分在扭变（twisting）之后成为 $S^+\otimes\mathcal{L}$ 的一个截面。我们已指出，表达式 $S^+\otimes\mathcal{L}$ 应理解为指称一个 $\mathrm{Spin}^c$ 结构的非正式写法。我们用符号 $\lambda$ 既表示该 $\mathrm{Spin}^c$ 结构本身，也表示它的第一 Chern 类（在假定 $H^2(X,\mathbb{Z})$ 中没有二阶挠元时，后者被定义为 $\lambda=S^+\otimes\mathcal{L}$ 的第一 Chern 类的一半）。

假定我们要对某个多项式 $P$ 与某个 $\mathrm{Spin}^c$ 结构 $\lambda$ 计算 $\langle P(a)\rangle_\lambda$。上面概述的高能论证表明，这个期望值集中在几何流形 $\mathcal{M}_\lambda$ 上——它是

% ===== 原文 p.279 =====
$\mathrm{Spin}^c$ 结构 $\lambda$ 的经典方程之解模空间的底流形。由于 $\exp(-S)$ 在 $\mathcal{M}_\lambda$ 上的取值为 $q^{\lambda^2}$，我们有
\[
\langle P(a)\rangle_\lambda=q^{\lambda^2}\int_{\mathcal{M}_\lambda}P(a).
\]
由于 $a$ 是 $\mathcal{M}_\lambda$ 上的一个二形式，这个积分等于
\[
q^{\lambda^2}c_{n_\lambda}\int_{\mathcal{M}_\lambda}a^{n_\lambda},
\]
其中 $n_\lambda$ 是 $\mathcal{M}_\lambda$ 的虚维数的一半，而 $c_{n_\lambda}$ 是 $P(a)$ 在 $a=0$ 处的 Taylor 级数中 $a^{n_\lambda}$ 的系数。当然，
\[
SW_X(\lambda)=\int_{\mathcal{M}_\lambda}a^{n_\lambda}
\]
按定义就是流形 $X$ 的 Seiberg--Witten 不变量在 $\mathrm{Spin}^c$ 结构 $\lambda$ 上的取值。（按约定，若 $n_\lambda$ 不是整数，则所有这些积分都解释为零。）于是我们已把期望值 $\langle P(a)\rangle_\lambda$ 求为：$P(a)$ 的 Taylor 级数中 $a^{n_\lambda}$ 的系数乘以 $q^{\lambda^2}SW_X(\lambda)$。对后面某些计算来说，把它改写成
\[
\langle P(a)\rangle_\lambda=q^{\lambda^2}SW_X(\lambda)\cdot\operatorname*{Res}_{a=0}\left(\frac{da}{a^{1+n_\lambda}}P(a)\right)
\]
是方便的。对所有 $\mathrm{Spin}^c$ 结构求和给出
\[
\langle P(a)\rangle=\sum_\lambda\langle P(a)\rangle_\lambda
=\sum_\lambda q^{\lambda^2}SW_X(\lambda)\cdot\operatorname*{Res}_{a=0}\left(\frac{da}{a^{1+n_\lambda}}P(a)\right).
\]
当我们也允许一个二维算符时，有类似的公式。由于这个算符只是测量第一 Chern 类，我们得到
\[
\begin{aligned}
\left\langle P(a)\cdot\exp\left(\alpha\int_\Sigma\mathcal{O}^{(2)}\right)\right\rangle
&=\sum_\lambda q^{\lambda^2}\cdot\exp\bigl(\alpha\langle c_1(\lambda),\Sigma\rangle\bigr)\langle P(a)\rangle_\lambda\\
&=\sum_\lambda q^{\lambda^2}\cdot\exp\bigl(\alpha\langle c_1(\lambda),\Sigma\rangle\bigr)SW_X(\lambda)\cdot
\operatorname*{Res}_{a=0}\left(\frac{da}{a^{1+n_\lambda}}P(a)\right).
\end{aligned}
\]

% [Tutor Note: 从「系数」改写成留数，不只是记号上的方便：$\operatorname*{Res}_{a=0}\bigl(da\,P(a)/a^{1+n}\bigr)$ 按定义就是 $da\,P(a)/a^{1+n}$ 的 $a^{-1}$ 系数，即 $P$ 的 $a^n$ 系数。之所以用留数而不用「Taylor 系数」，是因为 $n_\lambda$ 随 $\lambda$ 变化，而留数写法对 $P$ 的具体形式、对 $n_\lambda$ 的符号都不需要分情形讨论；同时它把 p.281 中「在 $a=0$ 附近把被积函数写成 $\frac{da}{a}\times(\text{正则部分})$」这一步（即 $n_\lambda=0$ 的单型情形）变得自然。易错点：若漏掉 $da$ 而只写 $\operatorname*{Res}P(a)/a^{1+n}$，在 $n\ne0$ 时会得到完全不同的结果；此外 $\mathcal{O}^{(2)}$ 与 $\Sigma$ 的配对 $\langle c_1(\lambda),\Sigma\rangle$ 是整数值（在无挠假设下），而 $\alpha$ 是任意复数，二者不要与 $q^{\lambda^2}$ 的指数混同。]

现在让我们推广这个 $U(1)$ 规范理论的计算，允许更多保持拓扑不变性的项。弯曲流形 $X$ 上保持拓扑不变性的 Lagrange 量，其最一般的形式是
\[
\mathcal{L}(a,A,M)=\{Q,\cdot\}+\int K^4e(a)+\int g(a)\operatorname{Pf}(R_X)+\int h(a)L(R_X),
\]
其中 $R_X$ 是 $X$ 的 Riemann 曲率张量。拓扑不变性要求函数 $e,g,h$ 必须是全纯函数。当然，$\int\operatorname{Pf}(R_X)$ 给出 Euler 示性数的一个倍数，而 $\int L(R_X)$ 给出 $X$ 的符号差（signature）的一个倍数。

% [Tutor Note: 上式第一项写作 $\{Q,\cdot\}$ 是原文的简写，意思是「某个 $Q$-正合项」，即形如 $\{Q,\Lambda\}$ 的东西；它只影响非物理的（$Q$-正合的）部分，故不改变 $Q$-闭可观测量的期望值。真正的易错之处在后面：$\operatorname{Pf}(R_X)$ 是 Euler 形式，$\int\operatorname{Pf}(R_X)=\chi(X)$ 只在正确的归一化（$\frac{1}{(2\pi)^2}$ 之类）下成立，而 $\int L(R_X)=\frac13 p_1=\sigma(X)$ 也有归一化因子；本文把它们一律吸收进未知函数 $g,h$（即 $G=e^{-g},H=e^{-h}$）之中，所以读者不必、也不应试图从这两个公式里定出数值系数——这些系数最终被 p.281 所说的「五个普适常数」吞掉了。]

期望值现在是
\[
\langle P(a)\rangle=\sum_\lambda SW(\lambda)\operatorname*{Res}_{a=0}
\left\{\frac{da}{a^{1+n_\lambda}}P(a)\left(E(a)^{\lambda^2}G(a)^{\chi(X)}H(a)^{\sigma(X)}\right)\right\},
\]
其中 $G=e^{-g}$，$H=e^{-h}$，而
\[
E=\exp\left(-\frac{\partial^2e}{\partial a^2}\right).
\]
最后一个等式的理由是：由于 $K(a)=\psi$ 而 $K(\psi)=F$，我们有 $K(e(a))=e'(a)\psi$，以及
\[
K^2(e(a))=e''(a)\psi\wedge\psi+e'(a)K(\psi)=e''(a)\psi\wedge\psi+e'(a)F.
\]

% [Tutor Note: 因子 $E^{\lambda^2}$ 中的负号与指数位置都值得推敲。把 $e^{-S}$ 展开时，$\int K^4e(a)$ 里的 $e''(a)F\wedge F$ 贡献 $e^{-e''\!(a)\int F\wedge F}$；而按本讲约定，$\mathrm{Spin}^c$ 结构 $\lambda$ 上的瞬子数为 $\lambda^2$（这正是上文 $\exp(-S)|_{\mathcal{M}_\lambda}=q^{\lambda^2}$ 的意思），故该因子恰为 $\exp(-e''\lambda^2)=E^{\lambda^2}$，其中 $E=\exp(-\partial^2e/\partial a^2)$。若把 $F\wedge F$ 归一化为 $4\pi^2k$ 而非 $8\pi^2k$，则 $E$ 应改为 $\exp(-2\partial^2e/\partial a^2)$，最终 $2^{1+\frac14(7\chi+11\sigma)}$ 中的系数也随之改变。另请注意 $E,G,H$ 都是 $a$（等价地 $u$）的函数，所以它们与留数中的 $P(a)$ 谁先谁后不影响结果，但与 $q^{\lambda^2}$ 的相对位置（后者是常数）不要写错。]

% ===== 原文 p.280 =====
继续求导，我们发现
\[
K^4(e(a))=e''(a)F\wedge F+\psi\wedge(\cdots).
\]
在我们的情形中，一切含 $\psi$ 的项都消失，因为 $b_1(X)=0$。若 $\tau$ 为常数，则因子 $E(a)^{\lambda^2}$ 就约化为 $q^{\lambda^2}$，如上所写。

% [Tutor Note: 为何 $b_1(X)=0$ 就能丢掉 $\psi\wedge(\cdots)$：扭变之后的费米子 $\psi$（即 $\psi_\mu dx^\mu$）在低能有效理论中只剩零模，而 $\psi$ 是 $1$ 形式取值于 $\mathfrak{u}(1)$ 的场，其零模即 $H^1(X;\mathbb{R})$ 中的元素，故 $b_1(X)=0$ 时 $\psi\equiv0$，一切含 $\psi$ 的项消失。这也是前面 p.278 处同一条理由的重复使用；若 $b_1\ne0$，则 $K^4u$ 的插入还会给出与 $\psi$ 零模有关的额外贡献，$u$ 平面上的被积函数就不再只是 $E,G,H$ 的乘积。]

\subsection{$SU(2)$ 的情形}
% 原文 §19.4.3：The case of $SU(2)$

现在我们准备回到纯 $SU(2)$ 理论，并用我们刚做过的 $U(1)$ 计算来理解低能有效计算。假定我们固定 $\alpha\in\mathbb{C}$ 以及一个类 $\Sigma\in H^2(X;\mathbb{C})$。我们要计算
\[
\left\langle\exp\left(\alpha u+\beta\int_\Sigma\mathcal{O}^{(2)}\right)\right\rangle_{SU(2)}.
\]
答案的一般形式将是
\[
\left\langle\exp\left(\alpha u+\beta\int_\Sigma\mathcal{O}^{(2)}\right)\right\rangle_{u=1}
+\left\langle\exp\left(\alpha u+\beta\int_\Sigma\mathcal{O}^{(2)}\right)\right\rangle_{u=-1}
+\int_{u\text{-plane}}\cdots.
\]
$u$ 平面积分的被积函数（它实际上是一个只依赖于初等拓扑数据的、非常复杂的表达式）就简记为 $\cdots$。

让我们看看在低能有效理论中得到什么算符。利用 $K$ 与取低能极限这一步对易这一事实，再加上上面关于 $K^2(e(a))$ 的计算，并取 $e(a)=u$，我们有
\[
\int_\Sigma\mathcal{O}^{(2)}\mapsto\int_\Sigma K^2u=\int_\Sigma\frac{\partial u}{\partial a}F+\int_\Sigma\frac{\partial^2u}{\partial a^2}\psi\wedge\psi.
\]
这里我们假定：在 $u$ 平面的某个局部区域内，已经选好了 $a$ 的一个好的取法；公式在该区域内、对于该取法成立。例如在 $u=\infty$ 附近我们可以用通常的 $a$，但在 $u=1$ 附近我们必须改用另一个 $a$。若 $b_1(X)=0$，则可以丢掉第二项，于是有
\[
\int_\Sigma\mathcal{O}^{(2)}\mapsto\int_\Sigma\frac{\partial u}{\partial a}F.
\]
现在利用上面得出的二维算符乘法规则，我们看到
\[
\exp\left(\alpha u+\beta\int_\Sigma\mathcal{O}^{(2)}\right)
\mapsto\exp\left(\alpha u+\beta\frac{\partial u}{\partial a}\int_\Sigma F+\beta^2\Sigma^2T(u)\right).
\]

% [Tutor Note: 这里「必须改用另一个 $a$」是作者一带而过的关键跳步，它正是 Seiberg--Witten 解的全纯结构：$u$ 与 $a$ 的关系由 SW 曲线（椭圆曲线族）给出，$a(u)$ 是多值的，且在 $u=\pm1$（磁单极子无质量点）附近，$a$ 的好的局部坐标需要用对偶坐标 $a_D$ 而不是微扰的 $a$。换句话说，$\frac{\partial u}{\partial a}$ 只在某个局部全纯坐标卡内才有定义，跨卡时按 $SL(2,\mathbb{Z})$ 的单值变换换算——这就是为什么作者特意声明「公式在该区域内、对该取法成立」。而 $\beta^2\Sigma^2T(u)$ 这一项来自 $\exp(\beta\int_\Sigma\mathcal{O}^{(2)})$ 展开的二次项与 p.278 的接触项公式：两条闭链都取 $\Sigma$ 时，相交配对 $(\Sigma_1\cap\Sigma_2)$ 就变成自交数 $\Sigma\cdot\Sigma=\Sigma^2$，而 $T(u)$ 的系数已按此重新定义。注意此处用的是自交数 $\Sigma^2$（可为负），不是 $\Sigma$ 与某个 $\mathrm{Spin}^c$ 结构 $\lambda$ 的配对 $(\Sigma,\lambda)$——后者出现在 p.281 的公式里，两者不可混用。]

现在让我们计算 Donaldson 不变量中 $u=1$ 的贡献。我们这里所用的 $a$ 是那个在 $u=1$ 处无质量的磁单极子场所产生的 $a$。我们平移这个坐标，使得 $u=1$ 对应于 $a=0$。低能有效理论是一个 $U(1)$ 理论，并且是上一节所讨论理论的特例，带有某些未知的普适函数 $e,g,h$（等价地，未知函数 $E,G,H$）。

% ===== 原文 p.281 =====
于是我们有
\[
\begin{aligned}
&\left\langle\exp\left(\alpha u+\beta\frac{\partial u}{\partial a}(\Sigma,\lambda)+\beta^2\Sigma^2T(u)\right)\right\rangle\\
&\quad=2\sum_\lambda SW(\lambda)\operatorname*{Res}_{a=0}\Biggl\{\frac{da}{a^{1+n_\lambda}}
\exp\Bigl(\alpha u+\beta\frac{\partial u}{\partial a}(\Sigma,\lambda)\\
&\qquad\qquad+\beta^2\Sigma^2T(u)\Bigr)E^{\lambda^2}H^{\sigma(X)}G^{\chi(X)}\Biggr\}.
\end{aligned}
\]
右端出现的乘法因子 $2$ 是出于如下理由而引入的。规范变换群通过 $SU(2)$ 丛的自同构起作用，但它在联络空间上的作用不是忠实的。存在一个二阶的中心群，它平凡地作用。在规范理论的路径积分中，必须除以规范变换群的体积来做归一化。这个群比忠实作用的那一群大一倍，这一事实使得在把路径积分与常规的拓扑公式相比较时引入一个因子 $2$——后者是通过在模空间上积分得到的，并不除以 $2$。

% [Tutor Note: 这个「因子 2」是 Donaldson 理论与 Seiberg--Witten 理论对账时最经典的一个坑。技术上讲：$SU(2)$ 主丛的自同构群与 $SO(3)$ 主丛的自同构群相差 $\mathbb{Z}_2$（即 $-1\in SU(2)$ 生成的、在伴随表示上平凡作用的中心元），故 $\text{Aut}(P_{SU(2)})=\text{Aut}(P_{SO(3)})\times\mathbb{Z}_2$；而 Donaldson 不变量的模空间是按 $SO(3)$ 丛（等价地按 $w_2\ne0$ 的 $SU(2)$ 丛）定义的，因此场论路径积分（除以整个 $SU(2)$ 规范群体积）会比拓扑定义多除以 $2$、即多出一个因子 $2$。原文此处说「在比较中引入因子 2」，方向与此一致：$2\times$（拓扑答案）。另请注意同一段中「Let the virtual dimension of the SW moduli space by $2n_\lambda$」的 $by$ 显系 $be$ 之误。]

令 SW 模空间的虚维数为 $2n_\lambda$。由指标定理，$n_\lambda$ 是 $\lambda^2,\chi(X),\sigma(X)$ 的线性函数；我们不需要它的精确公式。

% [Tutor Note: 作者说「不需要精确公式」，但为便于核对，这里给出标准结果：SW 模空间的虚维数为 $\frac{\lambda^2-(2\chi(X)+3\sigma(X))}{4}$（见 Witten, Monopoles and four-manifolds, Math. Res. Lett. 1 (1994)），故 $n_\lambda=\frac{\lambda^2-2\chi-3\sigma}{8}$。用它可验证后文 p.281 的断言：在单型（$n_\lambda=0$）情形下 $\lambda^2$ 可以由 $\chi,\sigma$ 线性表出，即 $\lambda^2=2\chi+3\sigma$，于是 $E^{\lambda^2}$ 可并入 $\tilde G^{\chi}\tilde H^{\sigma}$ 之中——这正是 $\tilde G,\tilde H$ 被称为「复数」而不含 $a$ 的原因。]

让我们假定流形 $X$ 是（在 SW 理论意义下）单型的（of simple type）。按定义，这意味着所有不变量 $SW_X(\lambda)$ 除 $n_\lambda=0$ 的那些之外全为零。所有已知的、$b_2^+>1$ 的单连通四维流形都具有这一性质。于是我们的表达式简化为
\[
2\sum_\lambda SW_X(\lambda)\operatorname*{Res}_{a=0}\frac{da}{a}\exp(\cdots),
\]
其中未写出的表达式是 $a$ 的正则表达式。我们可以在 $u=1$ 处求值，得到
\[
2\sum_\lambda SW_X(\lambda)\exp\left(\alpha u\big|_{u=1}+(\Sigma,\lambda)\frac{du}{da}\Big|_{u=1}+\Sigma^2T(1)\right)
\tilde G^{\chi(X)}\tilde H^{\sigma(X)}.
\]
这里 $\tilde G$ 与 $\tilde H$ 是由 $E,G,H$ 令 $u=1$ 而得到的复数，其间还用了指标定理（回忆 $n_\lambda=0$）把 $\lambda^2$ 表为 $\chi(X),\sigma(X)$ 的线性组合。为了完全确定答案，我们必须求出
\[
u(1),\qquad\frac{\partial u}{\partial a}\Big|_{u=1},\qquad\tilde G,\qquad\tilde H,\qquad T(1).
\]
当然，我们还有 $u=-1$ 处的贡献，它与这一个相差一个对合 $u\mapsto-u$。

这类计算在 $b_2^+(X)>1$ 时给出 Donaldson 不变量的完整求值，因为在这种情况下（远离奇点）不存在对 $u$ 平面的连续积分。该公式表明：当 $b_2^+(X)>1$ 时，$X$ 的 Donaldson 多项式不变量完全由 $\mathrm{Spin}^c$ 结构上的函数 $SW_X(\lambda)$ 决定。这个显式关系涉及五个普适常数，它们已在上一段中列出。幸运的是，由于 Donaldson 不变量已在许多情形下被计算过，这些常数可以被确定下来。事实上，这可以用许多种方式做到，而所得结果总是彼此一致。

% [Tutor Note: 「五个普适常数」即上式的 $u(1),\ \partial u/\partial a|_{u=1},\ \tilde G,\ \tilde H,\ T(1)$（$u=-1$ 处的相应量由 $u\mapsto-u$ 对合给出，不是独立常数）。要独立地确定它们，可用的「实例」必须满足 $b_2^+>1$ 且为单型，例如 $K3$：它 $b_2^+=3$，由 $n_\lambda=\frac{\lambda^2-2\chi-3\sigma}{8}$ 得 $n_\lambda=0$ 仅当 $\lambda=0$，而 $SW_{K3}(0)=1$，故代入本页公式后左端与右端都只剩一项，可直接解出常数之间的关系；$S^2\times S^2$、$\mathbb{CP}^2\#\overline{\mathbb{CP}}^2$ 一类 $b_2^+=1$ 的流形因 $u$ 平面还有连续贡献而不能直接使用。另一条独立检验是与 Kronheimer--Mrowka 关于单型流形的结构定理比对：其 Donaldson 级数形如 $D_X=e^{u^2/2}\sum_\lambda SW_X(\lambda)e^{c_1(\lambda)}$，而本页公式的 $\exp$ 中恰含二次项 $\Sigma^2T(1)$，这正是 $e^{u^2/2}$ 的来源。]

% ===== 原文 p.282 =====
最终结果是：对于 $b_2^+(X)>1$ 且属于 Seiberg--Witten 单型的流形 $X$，我们有
\[
\begin{aligned}
&\left\langle\exp\left(\sum_a\beta_a\int_{\Sigma_a}\mathcal{O}^{(2)}+\alpha\mathcal{O}^{(0)}\right)\right\rangle\\
&\quad=2^{1+\frac14(7\chi+11\sigma)}\Biggl(\exp\left(\frac{v^2}{2}+2\alpha\right)\sum_\lambda SW_X(\lambda)e^{v\cdot c_1(\lambda)}\\
&\qquad\qquad+i^{\Delta}\exp\left(-\frac{v^2}{2}-2\alpha\right)\sum_\lambda SW_X(\lambda)e^{-iv\cdot c_1(\lambda)}\Biggr),
\end{aligned}
\]
其中 $v=\sum_a\beta_a[\Sigma_a]$，而 $\Delta=(\chi(X)+\sigma(X))/4$。右端两项中的第一项是 $u=1$ 处的贡献，第二项是 $u=-1$ 处的贡献。后者可以借助对合 $u\mapsto-u$ 由前者得到。

% [Tutor Note: 值得核对一下指数里那些系数。（i）$\Delta=(\chi+\sigma)/4$ 在 $b_1=0$ 时化为 $(1+b_2^+)/2$：因为 $\chi=2+b_2^++b_2^-$、$\sigma=b_2^+-b_2^-$，故 $\chi+\sigma=2+2b_2^+$。于是当 $b_2^+$ 为偶数时 $\Delta$ 是半整数，$i^\Delta$ 便带一个 $i^{1/2}=e^{i\pi/4}$ 式的相位；作同样的代入可得 $1+\frac14(7\chi+11\sigma)=\frac92+\frac92b_2^+-b_2^-$，它在 $b_2^+$ 为偶数时同样是半整数（例如 $S^2\times S^2$：$b_2^\pm=2$，指数为 $11.5$，$\Delta=1.5$）。两处同时出现「半整数次幂」提示：这些系数本质上是某些平方根，其分支必须与拓扑定义的相位约定（$\mathrm{Spin}^c$ 结构的符号、$i^\Delta$ 的分支）一并固定；此处仅按扫描页字面转写，若要与 Moore--Witten 原文逐一核对，请特别留意 $7\chi+11\sigma$ 的系数与 $i^\Delta$ 的相位。（ii）作为一个纯算术检验：$\mathbb{CP}^2$ 有 $\chi=3,\sigma=1$，指数为 $1+\frac14(21+11)=9$，即 $2^9$，与 $\Delta=1$ 相容；但 $\mathbb{CP}^2$ 的 $b_2^+=1$，本公式对它并不直接适用。]

在指数中，我们在这里于 $\alpha$ 前面放了一个凭空塞进去的因子 $2$；在此之前的这类公式中，$\alpha$ 的系数一直是 $1$。这一差异的原因是：我们在整个讨论中作了一个隐含的假定，现在为了与 Donaldson 理论比较，必须把它修正过来。回忆一下，低能理论中有一个质量标度 $\Lambda^2$，磁单极子在 $u=\pm\Lambda^2$ 处变为无质量。今天我们取了 $\Lambda=1$，因而所用的单位制使磁单极子在 $u=\pm1$ 处变为无质量。但没有理由认为这个重整化方案恰好就是与拓扑相容的那一个。所以我们应当带一个额外的未知变量 $\pm\Lambda^2$ 来做计算，然后通过与实例中的拓扑计算相比较来确定 $\Lambda$。（或者，若臂力足够强，我们也可以尝试先验地确定：在与我们固定的重整化方案配套时，究竟哪一个重整化方案、等价地哪一个 $\Lambda$ 值，会与标准的拓扑定义相合。）结果发现，为使场论计算给出通常的 Donaldson 多项式答案，我们应当把奇点放在 $u=\pm2$ 处；这反映了如下事实：按标准的拓扑定义，Kronheimer 与 Mrowka 的单型（simple type）条件是 $u^2=4$，其根在 $u=\pm2$。把奇点处的 $u$ 取为 $\pm2$ 而不是 $\pm1$，正好解释了为什么在最终公式的指数中出现的是 $\pm2\alpha$ 而不是 $\pm\alpha$。

% [Tutor Note: 这段是理解「为何重新标度 $\Lambda$」的关键，值得补一句作者未明说的对应关系：场论的坐标 $u_{\rm FT}$ 与拓扑的坐标 $U$ 相差 $U=2u_{\rm FT}$（即 $\Lambda=2$ 而非 $1$），于是 $U^2=4\iff u_{\rm FT}^2=1$，而 $\alpha U=2\alpha u_{\rm FT}$ 正是 $\pm2\alpha$ 的来源；同时 $SW_X$ 的求和指标 $\lambda$ 不变（$\lambda$ 是整上同调类，不随 $u$ 的标度改变），所以只有 $\alpha$、$v^2$ 这些与 $u$ 同量纲的量被标度。注意 $v$ 的定义式里 $\beta_a$ 与 $[\Sigma_a]$ 都没有被重新标度，这意味着 $\beta$ 的归一化与 $\alpha$ 的归一化在拓扑定义下是彼此独立的两个约定——这正是核对最终公式时最容易混淆的地方。]

\subsection{参考文献与交叉引用}
% 原文：References and cross references

本材料的基本参考文献是

\begin{quote}
N. Seiberg and E. Witten, \emph{Electric-Magnetic duality, monopole condensation, and confinement in $N=2$ supersymmetric Yang--Mills theory}, Nucl. Phys. B \textbf{426} (1994), pp. 19--52.

N. Seiberg and E. Witten, \emph{Monopoles, Duality and Chiral Symmetry Breaking In $N=2$ Supersymmetric QCD}, Nucl. Phys. \textbf{B431} (1994) 484.

E. Witten, \emph{Monopoles and four-manifolds}, Math. Res. Lett. \textbf{1} (1994), pp. 769--796.
\end{quote}

带超多重态的 $\mathcal{N}=2$ 理论之解的例子，在第二篇文章中给出。对 Donaldson 理论的应用在第三篇文章中做出，其详细展开见

\begin{quote}
G. Moore and E. Witten, \emph{Integration over the $u$-plane in Donaldson theory}, Adv. Theor. Math. Phys. \textbf{1} (1997), pp. 298--387.
\end{quote}

特别地，在那篇文献中发展了 $b_2^+=1$ 的理论，并且证明了：若 $b_2^+(X)>1$，则 $u$ 平面（除奇点之外）没有贡献。

% ===== 原文 p.283 =====
把模空间的 Euler 示性数与数学方面的计算相比较，其说明见

\begin{quote}
C. Vafa and E. Witten, \emph{A strong coupling test of $S$-duality}, Nuclear Physics B, \textbf{431} (1994), pp. 3--77.
\end{quote}

$\beta$ 函数与渐近自由（asymptotic freedom）、红外自由（infrared freedom）之间的关系，在 [I-Renorm Groups, §§2,3] 中描述。为使超对称整体化（globalize）而扭变理论的做法，在第 14 讲 §3 中描述。

% [Tutor Note: [I-Renorm Groups, §§2,3] 是本丛书（Quantum Fields and Strings: A Course for Mathematicians）第 I 卷中关于重整化群的章节代码，其编号体系与第 II 卷的讲次（Lecture 1–19）不同，读者不要把它误当作本卷的某一讲；本卷中与之对应的是第 18 讲开头对 $\beta$ 函数与 $\tau$ 的重整化的讨论。]

% 本页（印刷 p.1401）其余部分为空白。
